% INTRODUCTION - COMPLETE
\section{Introduction}

Glycans are essential biological molecules involved in protein folding, cell signaling, immune recognition, and pathogen interaction \citep{varki2017}. Over 50\% of human proteins undergo glycosylation, making accurate glycan structure prediction crucial for understanding biological mechanisms and developing therapeutics \citep{helenius2004}. GlyTouCan, the international glycan structure repository, catalogs over 200,000 unique glycan structures \citep{tiemeyer2016}, highlighting the scale of structural diversity that researchers need to navigate.

AlphaFold 3 (AF3) represents a breakthrough in biomolecular structure prediction, achieving unprecedented accuracy for protein-ligand complexes including glycans \citep{abramson2024}. However, Huang et al.\ (2025) identified a critical limitation: standard input formats produce stereochemically incorrect glycan structures. The SMILES format yields only $\sim$62\% stereochemistry accuracy due to epimer confusion and anomeric inversion errors. The userCCD format improves to $\sim$82\% but still introduces linkage position errors. Only the CCD+bondedAtomPairs (BAP) format achieves high accuracy ($>$98\%) by explicitly specifying atom-to-atom bonds---yet requires 30--60 minutes of manual specification per structure.

We present GlycoSMILES2BAP, an automated pipeline that converts standard glycan notations directly to AF3-compatible CCD+BAP format. Our approach addresses three technical challenges: (1) mapping 28+ monosaccharide configurations to correct CCD codes, (2) tracking anomeric positions that differ between aldoses (C1) and ketoses such as sialic acids (C2), and (3) generating explicit atom-pair bond specifications for branched glycans.

%==============================================================================
% METHODS - NO DUPLICATES
%==============================================================================
\section{Methods}

\subsection{Pipeline Architecture}

GlycoSMILES2BAP operates through four sequential modules: Input Parsing, CCD Mapping, BAP Generation, and Output Formatting. The pipeline accepts IUPAC-condensed, WURCS, and GlycoCT input formats.

\subsection{CCD Mapper Module}

The CCD (Chemical Component Dictionary) provides standardized residue identifiers. Our mapper supports 28+ monosaccharide configurations. Key design decisions include: (1) case-insensitive matching, (2) anomeric position tracking (C2 for sialic acids, C1 for aldoses), and (3) ring oxygen positions (O4 for pentoses, O5 for hexoses, O6 for sialic acids).

\begin{table}[h]
\centering
\caption{Key CCD Mappings}
\begin{tabular}{lllll}
\toprule
Monosaccharide & Anomer & Config & CCD Code & Anomeric C \\
\midrule
GlcNAc & $\beta$ & D & NAG & C1 \\
Man & $\alpha$ & D & MAN & C1 \\
Gal & $\beta$ & D & GAL & C1 \\
Fuc & $\alpha$ & L & FUC & C1 \\
Neu5Ac & $\alpha$ & D & SIA & C2 \\
\bottomrule
\end{tabular}
\label{tab:ccd}
\end{table}

\subsection{BAP Generator Module}

The BAP generator produces explicit atom-pair bond specifications. Each glycosidic linkage is represented as: (donor\_residue, donor\_atom, acceptor\_residue, acceptor\_atom).

\noindent\textbf{Algorithm: BAP Generation}
\begin{enumerate}
\item For each linkage $l$ in glycan topology:
\item \quad Identify donor anomeric carbon (C1 for aldoses, C2 for ketoses)
\item \quad Identify acceptor oxygen position from linkage notation
\item \quad Generate atom pair tuple
\item Return list of all atom pairs
\end{enumerate}

\noindent The algorithm runs in $O(n)$ time where $n$ is the number of residues.

\subsection{PDB Structure Validation Method}

To validate stereochemistry handling, we compared GlycoSMILES2BAP output against known correct PDB structures. This approach was necessary because AF3 model parameters were not accessible. Four PDB structures were selected to test critical cases (Table~\ref{tab:pdb_cases}).

